% Owner: Yashpal (Member 1)
\subsection{Caches and the three kinds of miss}
A cache keeps recently used lines of memory close to the processor, so it relies on temporal
locality (the same word is used again soon) and spatial locality (its neighbours are used soon)
\cite{smith1982}. In a direct-mapped cache each line of memory can live in exactly one cache
line, chosen by the index bits of its address. Hill argues that direct mapping is attractive
because the hit path is short: there is only one tag to compare and no way to select
\cite{hill1988}. Hill and Smith divide misses into three classes
\cite{hill1989,hp_quant}:
\begin{itemize}
  \item \emph{compulsory} misses, the first reference to a line;
  \item \emph{capacity} misses, which a fully associative cache of the same size would also
        suffer, because the working set does not fit;
  \item \emph{conflict} misses, caused by two lines that map to the same set evicting each
        other even though the cache has room.
\end{itemize}
Conflict misses are the weak point of direct mapping. Jouppi showed that a small fully
associative victim cache removes many of them \cite{jouppi1990}. In a unified cache, code and
data can also conflict with \emph{each other}, which is the main effect we observe in
Section~\ref{sec:split}.

\subsection{Write policies and write buffers}
A write-through cache sends every store to the next level. A write-back cache keeps dirty lines
and writes them back only on eviction. On a write miss, a write-allocate cache fetches the line
first, while a no-write-allocate (write-around) cache sends the store to memory and leaves the
cache unchanged \cite{jouppi1993,hp_quant}. Jouppi separates write-hit from write-miss policies,
compares write-around and write-validate on write misses, and measures the traffic that
write-through and write-back caches send to the next level \cite{jouppi1993}. Write-through with
no-write-allocate is the simplest combination, because memory always holds the current data and
no line ever has to be written back, so we use it here. A write buffer lets the
processor continue as soon as the store has been accepted, but Skadron and Clark identify three
ways in which it can still stall the processor: the buffer is full, the buffer competes with a
cache miss for the next level, or a load needs data that is still in the buffer
\cite{skadron1997}. Our cache has a one-entry buffer, and all three cases appear in our
measurements.
